% !TEX root = ../main.tex
\corrigesChapitre{14}{Deuxième loi de Newton}
\corrigesSouples

\begin{blocLivre}{cor-14-1}
\begin{corrige}{1}{Où est le centre de masse ?}
\demarche{Le centre de masse est sur les éléments de symétrie et plus près des parties les plus massives.}
\begin{questions}
\item Le disque homogène et l'haltère symétrique ont un centre de symétrie : G est \res{\text{confondu avec ce centre}}.
\item La masse est plus concentrée du côté de la grosse boule : G s'en rapproche. Avec $d_1=\qty{10}{cm}$ et
$d_2=\qty{30}{cm}$ : $m_1\,d_1=3{,}0\times10=30$ et $m_2\,d_2=1{,}0\times30=30$ : \res{\text{l'égalité est vérifiée}}.
\item G est au centre de l'anneau, \res{\text{en dehors de la matière}}. On ne peut donc pas poser l'anneau sur une pointe
placée en G : il n'y a pas de matière à cet endroit.
\end{questions}
\end{corrige}
\end{blocLivre}
\begin{blocLivre}{cor-14-2}
\begin{corrige}{2}{Galiléen ou non ?}
\demarche{Un référentiel est galiléen si le principe d'inertie y est vérifié.}
\begin{questions}
\item a. \res{\text{Terrestre}} : la chute dure moins d'une seconde, la rotation de la Terre est négligeable.
b. \res{\text{Géocentrique}} : il est en translation quasi uniforme par rapport aux étoiles pendant la durée d'une
orbite. c. \res{\text{Héliocentrique}} : pour une planète, les étoiles lointaines fournissent la référence.
\item \res{\text{Non}} : la balle se met en mouvement alors qu'aucune force horizontale ne s'exerce sur elle (ses forces
se compensent) ; le principe d'inertie n'est pas vérifié.
\item Dans le référentiel terrestre, galiléen, la balle conserve sa vitesse : elle continue en ligne droite à vitesse
constante, tandis que le bus ralentit.
\end{questions}
\end{corrige}
\end{blocLivre}
\begin{blocLivre}{cor-14-3}
\begin{corrige}{3}{La pierre de curling}
\demarche{Le poids et la réaction normale se compensent ; seul le frottement agit horizontalement.}
\begin{questions}
\item Poids $\vec{P}$ (vertical, vers le bas), réaction normale $\vec{N}$ (verticale, vers le haut), frottement $\vec{f}$
(horizontal, opposé au mouvement).
\item Sur l'axe horizontal orienté dans le sens du mouvement : $-f=m\,a$, d'où
$a=-\dfrac{1{,}9}{19}=\res{\qty{-0,10}{m.s^{-2}}}$ : accélération \res{\text{opposée à la vitesse}}.
\item Mouvement \res{\text{rectiligne uniformément ralenti}} ; $v(t)=2{,}0-0{,}10\,t$ s'annule à \res{t=\qty{20}{s}}.
\item $d=2{,}0\times20-\frac12\times0{,}10\times20^2=\res{\qty{20}{m}}$.
\end{questions}
\schema[corrige]{rev-14-ex03-cor}
\legende{La vitesse indique le mouvement ; elle ne fait pas partie du bilan des forces.}
\end{corrige}
\end{blocLivre}
\begin{blocLivre}{cor-14-4}
\begin{corrige}{4}{Le pèse-personne dans l'ascenseur}
\demarche{Sur un axe vertical vers le haut : $N-m\,g=m\,a$, donc $N=m(g+a)$.}
\begin{questions}
\item Le poids $\vec{P}$ et la force $\vec{N}$ exercée par le pèse-personne, verticale vers le haut.
\item Référentiel terrestre, supposé galiléen : $N-m\,g=m\,a_z$.
\item Démarrage : $N=60\times(9{,}81+1{,}2)=\qty{661}{N}$, affichage \res{67~kg} ; vitesse constante : $N=m\,g$,
affichage \res{60~kg} ; arrêt : $N=60\times(9{,}81-1{,}2)=\qty{517}{N}$, affichage \res{53~kg}.
\item Au démarrage, elle se sent \res{\text{plus lourde}} ; à l'arrêt, \res{\text{plus légère}}.
\end{questions}
\end{corrige}
\end{blocLivre}
\begin{blocLivre}{cor-14-5}
\begin{corrige}{5}{Des systèmes en équilibre}
\demarche{Mouvement rectiligne uniforme ou immobilité : la somme des forces est nulle.}
\begin{questions}
\item a. Poids et frottement de l'air ; b. poids et tension de la chaîne. Dans les deux cas, le centre de masse est
immobile ou en mouvement rectiligne uniforme : d'après le principe d'inertie, \res{\text{les forces se compensent}}.
\item $f=m\,g=90\times9{,}81=\res{\qty{883}{N}}$ et $T=12\times9{,}81=\res{\qty{118}{N}}$.
\item \res{\text{Non}} : il est en mouvement rectiligne uniforme, et non immobile ; seul le lustre est à l'équilibre au
sens du cours.
\end{questions}
\end{corrige}
\end{blocLivre}
\begin{blocLivre}{cor-14-6}
\begin{corrige}{6}{Le skieur sur la pente}
\demarche{Projeter sur la pente (sens du mouvement) puis perpendiculairement à la pente.}
\begin{questions}
\item Poids $\vec{P}$ vertical, réaction normale $\vec{N}$ perpendiculaire à la pente, frottement $\vec{f}$ parallèle à
la pente, vers le haut (voir schéma).
\item Sur l'axe parallèle à la pente, vers le bas : $m\,g\sin\alpha-f=m\,a$, d'où
$a=g\sin\alpha-\dfrac{f}{m}=3{,}36-1{,}60=\res{\qty{1,76}{m.s^{-2}}}$.
\item $v=a\,t=\res{\qty{14}{m.s^{-1}}}$ et $d=\frac12\,a\,t^2=\res{\qty{56}{m}}$.
\item Perpendiculairement : $N-m\,g\cos\alpha=0$, d'où $N=75\times9{,}81\times\cos20^\circ=\res{\qty{691}{N}}$.
\end{questions}
\schema[corrige]{14-ex06}
\end{corrige}
\end{blocLivre}
\begin{blocLivre}{cor-14-7}
\begin{corrige}{7}{Du mouvement aux forces}
\demarche{Connaissant l'accélération, la deuxième loi donne la somme des forces : $\sum\vec{F}=m\,\vec{a}$.}
\begin{questions}
\item $a=\dfrac{\Delta v}{\Delta t}=\dfrac{27{,}8}{9{,}0}=\res{\qty{3,1}{m.s^{-2}}}$.
\item $\sum F=m\,a=1300\times3{,}09=\res{\qty{4,0e3}{N}}$, horizontale, vers l'avant : c'est la force de frottement exercée
par la route sur les roues motrices.
\item Mouvement circulaire uniforme : $\vec{a}$ centripète, $a=\dfrac{v^2}{R}=\dfrac{10^2}{25}=\qty{4,0}{m.s^{-2}}$ ; somme des
forces centripète, de valeur $m\,a=\res{\qty{5,2e3}{N}}$.
\item Le \res{\text{frottement latéral des pneus}} sur la route. Sur le verglas, il serait trop faible : la voiture
continuerait presque en ligne droite.
\end{questions}
\end{corrige}
\end{blocLivre}
\begin{blocLivre}{cor-14-8}
\begin{corrige}{8}{Le décollage d'une fusée}
\demarche{Sur un axe vertical vers le haut : $F-m\,g=m\,a$.}
\begin{questions}
\item Poids et poussée. Si la poussée est inférieure au poids, la somme des forces est dirigée vers le bas ou nulle :
la fusée reste posée sur son pas de tir.
\item $a=\dfrac{F}{m}-g=\dfrac{\num{1.3e7}}{\num{7.8e5}}-9{,}81=\res{\qty{6,9}{m.s^{-2}}}$.
\item $v=a\,t=\res{\qty{69}{m.s^{-1}}}$ et $h=\frac12\,a\,t^2=\res{\qty{3,4e2}{m}}$.
\item La fusée brûle son carburant : sa masse diminue rapidement. À poussée constante, $a=F/m-g$ augmente.
\end{questions}
\end{corrige}
\end{blocLivre}
\begin{blocLivre}{cor-14-9}
\begin{corrige}{9}{Une bille dans le glycérol}
\demarche{Deuxième loi sur un axe vertical vers le bas : $m\,g-\Pi-k\,v=m\,a$.}
\begin{questions}
\item Sur un axe vertical orienté vers le bas : $m\,a=m\,g-\Pi-k\,v$, d'où $\res{a=g-\dfrac{\Pi}{m}-\dfrac{k\,v}{m}}$.
\item Pour $v=0$ : $a=9{,}81-\dfrac{0{,}0133}{\num{8.4e-3}}=\res{\qty{8,2}{m.s^{-2}}}$.
\item La vitesse croît puis se stabilise vers \qty{0,41}{m.s^{-1}}. Quand $a=0$ :
$v_\mathrm{L}=\dfrac{m\,g-\Pi}{k}=\res{\qty{0,409}{m.s^{-1}}}$.
\item Le frottement du glycérol, très visqueux, est bien plus grand que celui de l'air pour une vitesse donnée, et la
masse de la bille est petite : le temps caractéristique $m/k$ ne vaut que \qty{0,05}{s}.
\end{questions}
\end{corrige}
\end{blocLivre}
\begin{blocLivre}{cor-14-10}
\begin{corrige}{10}{Prendre un virage sans déraper}
\demarche{En virage, la somme des forces est centripète : $f=m\,v^2/R$, qui ne doit pas dépasser $\mu\,m\,g$.}
\begin{questions}
\item Le \res{\text{frottement latéral}} de la route sur les pneus, dirigé \res{\text{vers le centre du virage}} ; le
poids et la réaction normale se compensent.
\item Selon $\un$ : $f=m\,\dfrac{v^2}{R}$.
\item $m\,\dfrac{v^2}{R}\leq\mu\,m\,g$ donne $v\leq\sqrt{\mu g R}$ : sur route sèche,
$v_{\max}=\sqrt{0{,}80\times9{,}81\times50}=\qty{19,8}{m.s^{-1}}$, soit \res{71~\text{km/h}} ; sur route mouillée,
\res{50~\text{km/h}}.
\item La force nécessaire et la force maximale sont toutes deux proportionnelles à la masse : \res{m\text{ se simplifie}}.
\end{questions}
\end{corrige}
\end{blocLivre}
\begin{blocLivre}{cor-14-11}
\begin{corrige}{11}{Deux patineurs se repoussent}
\demarche{Système isolé : le centre de masse, initialement immobile, le reste.}
\begin{questions}
\item Poids et réactions de la glace, verticaux, qui se compensent ; les forces de poussée sont \emph{intérieures} au
système. Donc $\vec{a}_\mathrm{G}=\vec{0}$ : G, immobile au départ, \res{\text{reste immobile}}.
\item $m_\mathrm{A}\,v_\mathrm{A}=m_\mathrm{B}\,v_\mathrm{B}$, d'où $v_\mathrm{A}=\dfrac{50\times2{,}4}{80}=\res{\qty{1,5}{m.s^{-1}}}$, en sens
opposé à B.
\item $d_\mathrm{A}=\qty{3,0}{m}$ et $d_\mathrm{B}=\qty{4,8}{m}$ ; avec l'origine en G :
$x_\mathrm{G}=\dfrac{80\times(-3{,}0)+50\times4{,}8}{130}=\res{0}$.
\end{questions}
\end{corrige}
\end{blocLivre}
\begin{blocLivre}{cor-14-12}
\begin{corrige}{12}{Le saut de Felix Baumgartner}
\demarche{Tant que le frottement est négligeable, $a\approx g$ ; à la vitesse maximale, $a=0$ et le frottement compense le poids.}
\begin{questions}
\item La force de frottement, proportionnelle à $\rho\,v^2$, est alors négligeable : la deuxième loi se réduit à
$m\,\vec{a}=m\,\vec{g}$, soit $\res{a\approx g}$.
\item $v=\qty{377}{m.s^{-1}}$ : $t=\dfrac{v}{g}=\res{\qty{39}{s}}$ et $h=\frac12\,g\,t^2=\res{\qty{7,3}{km}}$.
\item À la vitesse maximale, $\ddt{v}=0$ : $\res{a=0}$. Donc $\frac12\,\rho\,C_xS\,v^2=m\,g$ et
$\rho=\dfrac{2\,m\,g}{C_xS\,v^2}=\res{\qty{0,020}{kg.m^{-3}}}$, ce qui correspond à une altitude d'environ \res{29~\text{km}}.
\item En descendant, l'air devient plus dense : le frottement dépasse le poids, l'accélération est dirigée vers le haut
et la vitesse diminue.
\end{questions}
\end{corrige}
\end{blocLivre}
\begin{blocLivre}{cor-14-13}
\begin{corrige}{13}{L'accéléromètre du smartphone}
\demarche{La vitesse est l'aire sous la courbe $a_z(t)$ ; la force du plancher vaut $N=m(g+a_z)$.}
\begin{questions}
\item 0 à 2~s : uniformément accéléré vers le haut ; 2 à 10~s : uniforme ; 10 à 12~s : uniformément ralenti.
L'ascenseur \res{\text{monte}}.
\item $v_{\max}=0{,}80\times2{,}0=\res{\qty{1,6}{m.s^{-1}}}$ ; $h=1{,}6+1{,}6\times8{,}0+1{,}6=\res{\qty{16}{m}}$.
\item $N=m(g+a_z)$ : \res{\qty{690}{N}}, puis \res{\qty{638}{N}}, puis \res{\qty{586}{N}}.
\item Immobile, l'accéléromètre ne distingue pas la pesanteur d'une accélération ; en retranchant $g$, le logiciel affiche
directement l'accélération de l'ascenseur, nulle à l'arrêt.
\end{questions}
\end{corrige}
\end{blocLivre}
\begin{blocLivre}{cor-14-14}
\begin{corrige}{14}{Le téléski}
\demarche{Projeter sur la piste (sens de la montée) et perpendiculairement à la piste.}
\begin{questions}
\item Poids $\vec{P}$, réaction normale $\vec{N}$, frottement $\vec{f}$ (vers le bas de la piste) et force $\vec{T}$ de
la perche, inclinée de $\beta$ sur la piste.
\item Le mouvement est rectiligne uniforme : d'après le principe d'inertie, \res{\text{la somme des forces est nulle}}.
\item Selon la piste : $T\cos\beta-m\,g\sin\alpha-f=0$, d'où $T=\dfrac{177{,}7+60}{\cos25^\circ}=\res{\qty{262}{N}}$ ;
perpendiculairement : $N=m\,g\cos\alpha-T\sin\beta=\res{\qty{552}{N}}$.
\item $a=\dfrac{2{,}5}{2{,}0}=\qty{1,25}{m.s^{-2}}$ et $T\cos\beta-m\,g\sin\alpha-f=m\,a$ : $T=\res{\qty{359}{N}}$.
\end{questions}
\schema[corrige]{14-ex14}
\end{corrige}
\end{blocLivre}
\begin{blocLivre}{cor-14-15}
\begin{corrige}{15}{Le lancer du marteau}
\demarche{Avant le lâcher, la tension du câble fournit l'accélération centripète ; après, seul le poids agit.}
\begin{questions}
\item La boule concentre presque toute la masse : G est \res{\text{très proche du centre de la boule}}, légèrement
décalé vers le câble.
\item Mouvement circulaire uniforme : $\vec{a}$ centripète, $a=\dfrac{v^2}{R}=\dfrac{27^2}{1{,}9}=\res{\qty{384}{m.s^{-2}}}$ ;
le câble exerce $T\approx m\,a=7{,}26\times384=\res{\qty{2,8e3}{N}}$ (le poids, \qty{71}{N}, est négligeable devant).
\item Le poids du lanceur vaut \qty{1,08e3}{N} : la tension en vaut près de \res{2{,}6\text{ fois}}. Pour ne pas être
entraîné, le lanceur penche vers l'arrière : son propre poids compense en partie la traction du câble.
\item Seul le \res{\text{poids}} agit : $\vec{a}_\mathrm{G}=\vec{g}$, et la trajectoire de G est une
\res{\text{parabole}}, même si le marteau tourne sur lui-même.
\end{questions}
\end{corrige}
\end{blocLivre}
